\documentclass[12pt,a4paper]{article}
\usepackage[T1]{fontenc}
\usepackage{amsmath,amssymb,mathrsfs,tikz,times,pifont}
\usepackage{enumitem}
\usepackage{multicol}
\usepackage{lmodern}
\usetikzlibrary{trees}
\newcommand\circitem[1]{%
\tikz[baseline=(char.base)]{
\node[circle,draw=gray, fill=red!55,
minimum size=1.2em,inner sep=0] (char) {#1};}}
\newcommand\boxitem[1]{%
\tikz[baseline=(char.base)]{
\node[fill=cyan,
minimum size=1.2em,inner sep=0] (char) {#1};}}
\setlist[enumerate,1]{label=\protect\circitem{\arabic*}}
\setlist[enumerate,2]{label=\protect\boxitem{\alph*}}
%%%::::::by chnini ameur :::::::%%%
\everymath{\displaystyle}
\usepackage[left=1cm,right=1cm,top=1cm,bottom=1.7cm]{geometry}
\usepackage[colorlinks=true, linkcolor=blue, urlcolor=blue, citecolor=blue]{hyperref}
\usepackage{array,multirow}
\usepackage[most]{tcolorbox}
\usepackage{varwidth}
\usepackage{float} %pour utiliser l'option [H] qui force l'image à apparaître exactement à l'endroit où elle est placée dans le code.
\tcbuselibrary{skins,hooks}
\usetikzlibrary{patterns}
%%%::::::by chnini ameur :::::::%%%
\newtcolorbox{exa}[2][]{enhanced,breakable,before skip=2mm,after skip=5mm,
colback=yellow!20!white,colframe=black!20!blue,boxrule=0.5mm,
attach boxed title to top left ={xshift=0.6cm,yshift*=1mm-\tcboxedtitleheight},
fonttitle=\bfseries,
title={#2},#1,
% varwidth boxed title*=-3cm,
boxed title style={frame code={
\path[fill=tcbcolback!30!black]
([yshift=-1mm,xshift=-1mm]frame.north west)
arc[start angle=0,end angle=180,radius=1mm]
([yshift=-1mm,xshift=1mm]frame.north east)
arc[start angle=180,end angle=0,radius=1mm];
\path[left color=tcbcolback!60!black,right color = tcbcolback!60!black,
middle color = tcbcolback!80!black]
([xshift=-2mm]frame.north west) -- ([xshift=2mm]frame.north east)
[rounded corners=1mm]-- ([xshift=1mm,yshift=-1mm]frame.north east)
-- (frame.south east) -- (frame.south west)
-- ([xshift=-1mm,yshift=-1mm]frame.north west)
[sharp corners]-- cycle;
},interior engine=empty,
},interior style={top color=yellow!5}}
%%%%%%%%%%%%%%%%%%%%%%%
\usepackage{fancyhdr}
\usepackage{eso-pic}         % Pour ajouter des éléments en arrière-plan
% Commande pour ajouter du texte en arrière-plan
\usepackage{tkz-tab}
\AddToShipoutPicture{
    \AtTextCenter{%
        \makebox[0pt]{\rotatebox{80}{\textcolor[gray]{0.7}{\fontsize{5cm}{5cm}\selectfont PGB}}}
    }
}
\usepackage{lastpage}
\fancyhf{}
\pagestyle{fancy}
\renewcommand{\footrulewidth}{1pt}
\renewcommand{\headrulewidth}{0pt}
\renewcommand{\footruleskip}{10pt}
\fancyfoot[R]{
\color{blue}\ding{45}\ \textbf{2025}
}
\fancyfoot[L]{
\color{blue}\ding{45}\ \textbf{Prof:M. BA}
}
\cfoot{\bf
\thepage /
\pageref{LastPage}}
% Création du compteur pour les exercices
\newcounter{exercice}
\renewcommand{\theexercice}{\arabic{exercice}}  % Définit l'affichage du compteur en chiffres arabes

% Définir la commande \exo
\newcommand{\exo}{\refstepcounter{exercice}\textbf{Exercice \theexercice} }
\begin{document}
\renewcommand{\arraystretch}{1.5}
\renewcommand{\arrayrulewidth}{1.2pt}
\begin{tikzpicture}[overlay,remember picture]
    \node[draw=blue,line width=1.2pt,fill=purple,text=blue,inner sep=3mm,rounded corners,pattern=dots]at ([yshift=-2.5cm]current page.north) {\begingroup\setlength{\fboxsep}{0pt}\colorbox{white}{\begin{tabular}{|*1{>{\centering \arraybackslash}p{0.28\textwidth}} |*2{>{\centering \arraybackslash}p{0.2\textwidth}|} *1{>{\centering \arraybackslash}p{0.19\textwidth}|} }
                \hline
                \multicolumn{3}{|c|}{$\diamond$$\diamond$$\diamond$\ \textbf{Lycée de Dindéfélo}\ $\diamond$$\diamond$$\diamond$ } & \textbf{A.S. : 2025/2026}                                              \\ \hline
                \textbf{Matière: Mathématiques}                                                                                    & \textbf{Niveau : 4ème}\textbf{M2} & \textbf{Date: 09/01/2026} & \textbf{} \\ \hline
                \multicolumn{4}{|c|}{\parbox[c]{10cm}{\begin{center}
                                                                  \textbf{{\Large\sffamily Td Développement et Factorisation}}
                                                              \end{center}}}                                                                                                        \\ \hline
            \end{tabular}}\endgroup};
\end{tikzpicture}
\vspace{3cm}

\begin{multicols}{2}
\setlength{\columnseprule}{0.1mm} % La largeur de la ligne verticale entre les colonnes
\fbox{\textbf{\exo}} Développer et réduire les expressions suivantes :
\begin{multicols}{2}
\setlength{\columnseprule}{0.1mm} % La largeur de la ligne verticale entre les colonnes
\begin{enumerate}
    \item $3(x+5)$
    \item $-2(4x-1)$
    \item $(x+2)(x+7)$
    \item $(2x-3)(x+5)$
    \item $(x+4)^2$
    \item $(3x-1)^2$
\end{enumerate}
\end{multicols}

\fbox{\textbf{\exo}} Développer puis réduire
\begin{multicols}{2}
\setlength{\columnseprule}{0.1mm} % La largeur de la ligne verticale entre les colonnes
\begin{enumerate}
    \item $5(2x-3) - 2(x+4)$
    \item $(x+3)(x-3)$
    \item $(2x+1)(x+4)$
    \item $(x-5)^2 + 3x$
\end{enumerate}
\end{multicols}
\fbox{\textbf{\exo}} Développer ou factoriser

Indiquer s'il faut développer ou factoriser, puis effectuer le calcul.
\begin{multicols}{2}
\setlength{\columnseprule}{0.1mm} % La largeur de la ligne verticale entre les colonnes
\begin{enumerate}
    \item $(x+1)(x+9)$
    \item $6x + 12$
    \item $(2x-5)^2$
    \item $x^2 - 4$
\end{enumerate}
\end{multicols}

\fbox{\textbf{\exo}} Factoriser (mise en évidence)

Factoriser les expressions suivantes :
\begin{multicols}{2}
\setlength{\columnseprule}{0.1mm} % La largeur de la ligne verticale entre les colonnes
\begin{enumerate}
    \item $5x + 10$
    \item $3x^2 + 6x$
    \item $7x - 14$
    \item $x^2 + 4x$
    \item $2x^2 - 8x$
\end{enumerate}
\end{multicols}
\fbox{\textbf{\exo}} Factoriser à l'aide des identités remarquables
\begin{multicols}{2}
\setlength{\columnseprule}{0.1mm} % La largeur de la ligne verticale entre les colonnes
\begin{enumerate}
    \item $x^2 + 6x + 9$
    \item $x^2 - 16$
    \item $4x^2 - 12x + 9$
    \item $9x^2 - 25$
\end{enumerate}
\end{multicols}
\fbox{\textbf{\exo}} Factorisation avec facteur caché

Dans chaque expression, commencer par \textbf{ en évidence le facteur commun caché}.
\begin{enumerate}
\item $x(x+3) + 2(x+3)$
\item $3x(2x-1) - 5(2x-1)$
\item $(x-4)(2x+1) + x(2x+1)$
\item $5x(x-2) + 3(x-2)$
\item $(x+1)(x-3) - 2x(x-3)$
\item $4x(3x+5) - (3x+5)$
\end{enumerate}

\fbox{\textbf{\exo}} Facteurs communs cachés et identités remarquables

Factoriser complètement les expressions suivantes.
\begin{multicols}{2}
\setlength{\columnseprule}{0.1mm} % La largeur de la ligne verticale entre les colonnes

\begin{enumerate}
\item $(x+2)^2 + 3(x+2)$
\item $(x-5)^2 - 4(x-5)$
\item $2(x+1)^2 + 6(x+1)$
\item $(2x-3)^2 - (2x-3)$
\item $x(x-1)^2 + (x-1)^2$
\item $3(2x+1)^2 - (2x+1)$
\end{enumerate}
\end{multicols}
\fbox{\textbf{\exo}} Facteurs communs très cachés (niveau difficile)

Dans ces expressions, le facteur commun n'apparaît qu'après une première transformation.
\begin{enumerate}
\item $x^2 + 4x + 4 + 2x + 4$
\item $4x^2 - 4x + 1 - (2x - 1)$
\item $x^2 - 1 + x(x-1)$
\item $9x^2 + 6x + 1 - (3x+1)$
\item $x^2 - 6x + 9 + x - 3$
\item $4x^2 - 12x + 9 + 2x - 3$
\end{enumerate}

\fbox{\textbf{\exo}} piège ultime (annulations cachées)

\textbf{Attention :} dans chaque expression, certains termes s'annulent après réduction. Il faut extbf{réduire, puis factoriser}.
\begin{enumerate}
\item $(x+2)^2 - (x^2 + 4x)$
\item $(2x-3)^2 - (4x^2 - 12x + 9)$
\item $x(x-1) + (1-x)x$
\item $(x+5)(x-5) - (x^2 - 25)$
\item $3x(x+1) - x(3x+3)$
\item $(x-2)^2 + 4x - (x^2 - 4)$
\end{enumerate}

\fbox{\textbf{\exo}} : Problème

On considère le rectangle de longueur $(x+3)$ et de largeur $(x-1)$.
\begin{enumerate}
    \item Exprimer l'aire du rectangle sous forme développée.
    \item Factoriser l'expression obtenue.
\end{enumerate}
\end{multicols}
\end{document}